\documentclass[11pt,a4paper]{article}
\usepackage[T1]{fontenc}
\usepackage{isabelle,isabellesym}

% Keep pdfsetup last.
\usepackage{pdfsetup}

\urlstyle{rm}
\isabellestyle{it}
% Render informal proof comments, including underscores and Isabelle symbols.
\def\fregecommentsymbol#1>{\ensuremath{\csname #1\endcsname}}
\newcommand{\fregecommentbody}[1]{{\isastylecmt--- #1}\endgroup}
\renewcommand{\isamarkupcmt}{%
  \begingroup
  \catcode`\_=12
  \catcode`\^=12
  \catcode`\#=12
  \let\<\fregecommentsymbol
  \fregecommentbody}

\begin{document}

\title{Polynomial-Size Simulation of Frege Systems}
\author{Jakub Gonera}
\date{}
\maketitle

\begin{abstract}
  This entry formalizes a polynomial-size simulation result for Frege systems
  over possibly different finite connective alphabets. Assuming that both
  systems satisfy the formalized Frege-system conditions, it proves that one
  simulates the other with a polynomial bound on proof size. The result does
  not establish that the proof translation is computable in polynomial time.
\end{abstract}

\tableofcontents
\clearpage

\section*{Introduction}
The formalization follows the treatment of Frege-system simulation by
Reckhow~\cite{Reckhow1976}, Kraj\'i\v{c}ek~\cite{Krajicek1995}, and
Filmus~\cite{Filmus2010}.

\input{session}

\bibliographystyle{abbrv}
\bibliography{root}

\end{document}
